\% Integration point: place after Stage 8's local-rescue prerequisite and
\% before optional storage. This file deliberately contains no document wrapper.
\subsection{Join and prove Windows 11}

Start on a wired provisioning connection. Log in as the tested local rescue
administrator, set the permanent workstation name, and point the interface at
AD DNS. Do not join while the client is using public, gateway, or fallback
DNS.

\begin{shellblock}
Rename-Computer -NewName "<workstation-name>" -Restart
Get-DnsClientServerAddress -AddressFamily IPv4
Resolve-DnsName -Type SRV "\_ldap.\_tcp.dc.\_msdcs.<identity-domain>"
w32tm /query /source
Test-NetConnection <dc-fqdn> -Port 445
\end{shellblock}

\expectthis{The computer name is exact; every active domain-member interface
lists only approved AD DNS servers; the SRV answer names a DC; time has a
source; TCP 445 succeeds.}
\failurebranch{If ordinary names resolve but the SRV query fails, correct DHCP
or interface DNS. If the clock differs from the DC by five minutes or more,
correct time before touching the domain. Do not work around either fault with
a hosts-file entry.}

Join from \textit{Settings > System > About > Domain or workgroup > Change},
select \textit{Domain}, enter \placeholder{identity-domain}, and supply the
temporary delegated join credential. A domain administrator is not required
for routine joins. Restart when Windows asks.

\begin{shellblock}
Get-CimInstance Win32\_ComputerSystem |
  Select-Object Name,PartOfDomain,Domain
nltest /dsgetdc:<identity-domain>
Test-ComputerSecureChannel -Verbose
\end{shellblock}

\expectthis{\texttt{PartOfDomain} is true, \texttt{Domain} is exact,
\texttt{nltest} discovers a DC through DNS, and the secure channel reports
\texttt{True}. Record the DC name and test time.}

At the sign-in screen choose \textit{Other user}. Sign in once as the standard
test identity using \texttt{DOMAIN\textbackslash user} or the full user
principal name. Open a non-elevated PowerShell session:

\begin{shellblock}
whoami
whoami /groups
echo \$env:USERPROFILE
net localgroup Administrators
\end{shellblock}

\proofpair{The reported identity is the standard domain test user}{Its profile
is local, its SID is present, and it is absent from the local Administrators
group.}
\proofpair{Sign in as the approved daily administrator}{A deliberate elevation
prompt succeeds, but the reserved domain-administrator identity has not been
used for daily work.}
\proofpair{Disconnect every network interface and restart}{The previously used
domain identity signs in from its cached credential; network resources fail
without delaying or preventing the local desktop.}
\proofpair{Reconnect, disable the standard test identity, then attempt a fresh
online sign-in}{Connected authentication is denied. Record separately that an
already cached offline credential cannot be revoked in phase 1.}

\askbefore{Did the join use a delegated credential? Can local rescue still sign
in with the DC powered off? Is the standard identity non-administrative? Did
the test measure DNS answer, clock source, secure-channel state, and profile
path rather than infer success from a desktop appearing?}

\subsection{Join and prove Arch Linux}

Finish a full system update before the join; Arch does not support partial
upgrades. Install the declared identity client as one transaction:

\begin{shellblock}
sudo pacman -Syu --needed sssd samba krb5 bind
resolvectl status
timedatectl show -p NTPSynchronized --value
host -t SRV \_ldap.\_tcp.dc.\_msdcs.<identity-domain>
\end{shellblock}

\expectthis{AD DNS is active, time reports \texttt{yes}, the LDAP SRV answer
names a DC, and the packages came from enabled official Arch repositories.}
\failurebranch{Stop if discovery fails, time is not synchronized, or Network
Manager can replace AD DNS with a public resolver. A successful ping is not an
identity-plane test.}

Create a one-use delegated join identity or obtain a time-limited credential
from the encrypted secret channel. Generate the Samba configuration from
private values with \texttt{security = ads}, the exact
\texttt{workgroup = <netbios-name>} and
\texttt{realm = <kerberos-realm>}. Set the Kerberos method to
\texttt{secrets and keytab}. Generate the root-owned, mode-0600 SSSD
configuration with:

\begin{shellblock}
[sssd]
services = nss, pam
config\_file\_version = 2
domains = <identity-domain>

[domain/<identity-domain>]
id\_provider = ad
auth\_provider = ad
access\_provider = ad
ad\_domain = <identity-domain>
krb5\_realm = <kerberos-realm>
cache\_credentials = true
use\_fully\_qualified\_names = true
fallback\_homedir = /home/\%u
default\_shell = /bin/bash
ldap\_id\_mapping = true
\end{shellblock}

Review both rendered files before joining. Join interactively; never place the
password in a command, shell history, answer file, or repository.

\begin{shellblock}
sudo net ads join -U '<delegated-join-user>'
sudo net ads testjoin
sudo klist -k /etc/krb5.keytab
sudo test "\$(stat -c \%a /etc/sssd/sssd.conf)" = 600
sudo systemctl enable --now sssd
\end{shellblock}

\expectthis{\texttt{net ads testjoin} reports success, the machine keytab
contains the expected host principals, \texttt{/etc/sssd/sssd.conf} is
root-owned mode 0600, and \texttt{sssd.service} is active. Revoke or expire
the join credential now.}

Keep the naming rule explicit. Either retain fully qualified login names
(\texttt{user@domain}) or set a reviewed, site-wide short-name rule; never
allow different domains or local accounts to collapse onto the same visible
name. Confirm that \texttt{sss} appears on the \texttt{passwd}, \texttt{group},
and \texttt{shadow} lines in \texttt{/etc/nsswitch.conf}. Then verify lookups:

\begin{shellblock}
getent passwd '<standard-user>@<identity-domain>'
id '<standard-user>@<identity-domain>'
getent group '<approved-admin-group>@<identity-domain>'
sssctl user-checks '<standard-user>@<identity-domain>' -s login
\end{shellblock}

\expectthis{The user resolves once, has a stable numeric UID and GID, is absent
from the approved administrator group, and SSSD authorizes the login service.
Record the numeric IDs; optional SMB files must present compatible ownership.}

Local home creation is a PAM action, not a network-home mount. Before editing,
save the package-owned file's digest and a private evidence copy. Add this line
to the session stack used by the workstation login path, normally
\texttt{/etc/pam.d/system-login}, after the other required session setup:

\begin{shellblock}
session optional pam\_mkhomedir.so skel=/etc/skel umask=0077
\end{shellblock}

Review the complete PAM stack before rebooting. Do not add
\texttt{pam\_mkhomedir} to authentication or account sections, and do not
replace package-managed includes. Test first from a second TTY while the local
rescue session remains open.

\begin{shellblock}
getent passwd '<standard-user>@<identity-domain>'
findmnt --target '/home/<resolved-home>'
stat -c '\%U \%G \%a \%n' '/home/<resolved-home>'
sudo -l -U '<standard-user>@<identity-domain>'
\end{shellblock}

\proofpair{First online login}{Creates exactly the home path returned by
\texttt{getent}; it is a local filesystem path, owned by the resolved UID/GID,
and begins with restrictive permissions.}
\proofpair{Standard-user authorization}{The user receives no workstation
administration through \texttt{sudo}.}
\proofpair{Approved administrator authorization}{Only the reviewed domain
group receives the intended narrow \texttt{sudo} policy; a real elevation is
tested and logged.}
\proofpair{Offline restart after one successful online login}{The cached
identity signs in, the same numeric UID and GID are returned, and the existing
local home opens even when DNS, LDAP, Kerberos, and SMB are unavailable.}
\proofpair{Local rescue with the DC unavailable}{Login and \texttt{sudo}
succeed without SSSD. If they do not, restore the PAM/NSS configuration from
the still-open rescue session.}

\begin{gate}{UID, time, and storage proof}
Record UID, primary GID, supplementary groups, client time, DC time, and the
SMB file owner observed from both operating systems. A name that merely looks
correct is not proof. NFS remains disabled until numeric-ID mapping is designed
and tested; phase 1 uses SMB and a local home.
\end{gate}

\askbefore{What exact login spelling is canonical? Which AD group alone grants
local administration? Does the home exist and open with storage powered off?
Do UID and GID remain unchanged across two cold boots? Is clock offset below
the Kerberos limit? Does disabling an account deny a new connected login while
the documented offline-cache limitation remains understood?}
